\section{Тема 3. Циклические и итерационные алгоритмы}

\subsection{Постановка задачи (вариант~4)}

Прямоугольник на плоскости $[x_1,x_2]\times[y_1,y_2]$ задаётся
четырьмя числами (его габаритами). Последовательно вводятся
габариты прямоугольников; в процессе ввода находить площадь их
пересечения, не запоминая всех введённых прямоугольников.

\subsection{Математическое обоснование}

Пересечение двух прямоугольников
$[x_1,x_2]{\times}[y_1,y_2]$ и $[x_1',x_2']{\times}[y_1',y_2']$
само является прямоугольником
\[
  [\max(x_1,x_1'),\,\min(x_2,x_2')] \times
  [\max(y_1,y_1'),\,\min(y_2,y_2')],
\]
если границы по обеим осям не «перевёрнуты» (левая меньше правой,
нижняя меньше верхней); иначе пересечение пусто.

Поскольку операция пересечения ассоциативна, текущий результат
$R_i$ можно обновлять по мере ввода $i$-го прямоугольника:
\[
  R_i = R_{i-1} \cap P_i,
\]
не сохраняя предыдущие прямоугольники в массиве.

\subsection{Блок-схема алгоритма}

\begin{center}
\begin{tikzpicture}[
  node distance=10mm, start chain=going below,
  nodes={draw=black, top color=white, bottom color=lightgray!50,
         minimum width=5.6cm, align=center, font=\small},
  arrow/.style={->, >=Stealth, thick},
]
  \node[draw, rounded corners, on chain] {Start};
  \node[shape=trapezium, trapezium left angle=70,
        trapezium right angle=110, on chain, draw, join=by arrow]
       {$n$,\ $R \leftarrow P_0$};
  \node[shape=chamfered rectangle, chamfered rectangle xsep=10pt,
        on chain, draw, join=by arrow]
       (loop) {$i := 1,\ n{-}1$};
  \begin{scope}[local bounding box=sci]
    \node[shape=trapezium, trapezium left angle=70,
          trapezium right angle=110, on chain, draw, join=by arrow]
         {$P_i$};
    \node[shape=rectangle, on chain, draw, join=by arrow]
         {$R := R \cap P_i$};
    \coordinate[on chain, join=by arrow] (lend);
  \end{scope}
  \draw[arrow] (sci.south) -| ($(sci.west)-(0.9,0)$) |- (loop);
  \coordinate[on chain] (sp);
  \draw[arrow] (loop) -| ($(sci.east)+(0.9,0)$) |- (sp);
  \node[shape=rectangle, on chain, draw, join=by arrow]
       {$S := (R_{x2}{-}R_{x1})(R_{y2}{-}R_{y1})$};
  \node[shape=trapezium, trapezium left angle=70,
        trapezium right angle=110, on chain, draw, join=by arrow]
       {$S$};
  \node[draw, rounded corners, on chain, join=by arrow] {End};
\end{tikzpicture}
\end{center}

\subsection{Текст программы}

\lstinputlisting[style=Cstyle]{../src/t03.c}

\subsection{Тестовые данные}

\begin{center}
\begin{tabular}{l|l}
\toprule
Прямоугольники & Площадь пересечения \\
\midrule
$[0,10]{\times}[0,10]$, $[5,15]{\times}[5,15]$ & 25 \\
$[0,5]{\times}[0,5]$, $[10,15]{\times}[10,15]$ & 0 (не пересекаются) \\
\bottomrule
\end{tabular}
\end{center}

\textbf{Результат выполнения программы} (первый тест):
\begin{lstlisting}[language={},basicstyle=\ttfamily\small,frame=single]
Введите количество прямоугольников: 2
Первый прямоугольник (x1 y1 x2 y2): 0 0 10 10
Прямоугольник 2 (x1 y1 x2 y2): 5 5 15 15
  Пересечение: [5.00,10.00]x[5.00,10.00]
Площадь итогового пересечения: 25.0000
\end{lstlisting}
